\documentclass[a4paper]{article}
\usepackage[fontset=fandol]{ctex}
\usepackage{amsmath, amssymb}
\usepackage{graphicx}
\usepackage[margin=2.5cm]{geometry}
\usepackage[most]{tcolorbox}
\usepackage{etoolbox}
\usepackage{listings}
\usepackage{hyperref}
\usepackage{booktabs}
\usepackage{subcaption}
\usepackage{float}
\usepackage{tikz}
\IfFileExists{pgfplots.sty}{\usepackage{pgfplots}\pgfplotsset{compat=1.18}}{}
\newtcolorbox{knowledgebox}[1]{enhanced, colback=blue!5!white, colframe=blue!75!black, colbacktitle=blue!75!black, coltitle=white, fonttitle=\bfseries, title=#1, attach boxed title to top left={yshift=-2mm, xshift=2mm}, boxrule=1pt, sharp corners}
\newtcolorbox{importantbox}[1]{enhanced, colback=yellow!10!white, colframe=yellow!80!black, colbacktitle=yellow!80!black, coltitle=black, fonttitle=\bfseries, title=#1, sharp corners}
\newtcolorbox{warningbox}[1]{enhanced, colback=red!5!white, colframe=red!75!black, colbacktitle=red!75!black, coltitle=white, fonttitle=\bfseries, title=#1, sharp corners}
\lstset{language=Python,basicstyle=\ttfamily\small,keywordstyle=\color{blue},stringstyle=\color{red!60!black},commentstyle=\color{green!60!black},breaklines=true,frame=single,numbers=left,numberstyle=\tiny\color{gray},captionpos=b,extendedchars=false}
\newcommand{\notetitle}{End-to-End ASR: Attention}
\newcommand{\noteauthors}{Codex 基于课程页、公开视频字幕、coverage ledger 与补充官方材料整理}
\newcommand{\notedate}{\today}
\newcommand{\videochannel}{WAVLab}
\newcommand{\videopublishdate}{2023-11-06}
\newcommand{\videoduration}{unknown}
\newcommand{\videourl}{https://www.youtube.com/watch?v=1FsamWmRO7Q}
\newcommand{\videocoverpath}{cover.jpg}
\begin{document}
\begin{titlepage}
\centering
{\Large 课程笔记\par}
\vspace{1.2cm}
{\huge\bfseries \notetitle\par}
\vspace{0.8cm}
{\large \noteauthors\par}
\vspace{0.3cm}
{\large \notedate\par}
\vspace{1.2cm}
\ifdefempty{\videocoverpath}{{\small 公开视频封面不可用。\par}}{\includegraphics[width=0.82\textwidth,height=0.45\textheight,keepaspectratio]{\videocoverpath}\par}
\vfill
\begin{tcolorbox}[width=0.9\textwidth, colback=black!2!white, colframe=black!60, sharp corners]
\textbf{视频作者/频道}：\videochannel\par
\textbf{发布日期}：\videopublishdate\par
\textbf{视频时长}：\videoduration\par
\textbf{视频链接}：\href{\videourl}{\nolinkurl{\videourl}}
\end{tcolorbox}
\end{titlepage}
\tableofcontents
\newpage
\section{Attention-based seq2seq overview}
根据课程页主题、公开视频字幕以及本讲的 coverage ledger，Attention-based seq2seq overview 这一部分主要围绕 End-to-End ASR: Attention | End-to-End ASR: Attention；attention-based ASR 的总览、encoder-decoder 复习，以及该方法相对前序模型的定位。 展开。
写作上保持 coverage-first 原则：先把老师在这一阶段真正讲过的问题设定、模型部件、训练或推理逻辑讲清，再压缩成适合复习的结构。

\begin{figure}[H]
\centering
\includegraphics[width=0.92\textwidth]{figures/coverage_map.png}
\caption{本讲 coverage map：按 coverage\_units.jsonl 归纳出的核心主题与章节映射。}
\end{figure}

\begin{figure}[H]
\centering
\includegraphics[width=0.92\textwidth]{figures/timeline.png}
\caption{本讲 timeline：根据字幕时间范围归纳的主题推进顺序。}
\end{figure}

\subsection{End-to-End ASR: Attention}
从证据位置看，这一部分对应的课程证据区间是 2023-11-06。课程在这里强调的核心内容可以概括为：End-to-End ASR: Attention | End-to-End ASR: Attention。
与单调模型不同，attention encoder-decoder 更强调全局条件化能力，但也因此引入了对齐自由度和部署层面的权衡。

在本章结构中，这一单元被并入“Attention-based seq2seq overview”，目的是把同一阶段的定义、机制和权衡集中讲清。

\subsection{attention-based ASR 的总览、encoder-decoder 复习，以及该方法相对前序模型的定位}
从证据位置看，这一部分对应的课程证据区间是 00:00:00,480--00:07:26,720。课程在这里强调的核心内容可以概括为：attention-based ASR 的总览、encoder-decoder 复习，以及该方法相对前序模型的定位。。
与单调模型不同，attention encoder-decoder 更强调全局条件化能力，但也因此引入了对齐自由度和部署层面的权衡。

在本章结构中，这一单元被并入“Attention-based seq2seq overview”，目的是把同一阶段的定义、机制和权衡集中讲清。

\subsection{本章小结}
Attention-based seq2seq overview 这一部分把 End-to-End ASR: Attention 中相邻的教学单元串成了一条连续叙事：既交代了这一阶段要解决的问题，也交代了方法、推理或工程权衡为什么要这样组织。

\section{Attention computation and context vectors}
根据课程页主题、公开视频字幕以及本讲的 coverage ledger，Attention computation and context vectors 这一部分主要围绕 attention 权重、context vector 与 weighted sum 如何把 encoder 信息接到 decoder。 展开。
写作上保持 coverage-first 原则：先把老师在这一阶段真正讲过的问题设定、模型部件、训练或推理逻辑讲清，再压缩成适合复习的结构。

\subsection{attention 权重、context vector 与 weighted sum 如何把 encoder 信息接到 decoder}
从证据位置看，这一部分对应的课程证据区间是 00:07:23,879--00:14:55,639。课程在这里强调的核心内容可以概括为：attention 权重、context vector 与 weighted sum 如何把 encoder 信息接到 decoder。。
与单调模型不同，attention encoder-decoder 更强调全局条件化能力，但也因此引入了对齐自由度和部署层面的权衡。

在本章结构中，这一单元被并入“Attention computation and context vectors”，目的是把同一阶段的定义、机制和权衡集中讲清。

\subsection{本章小结}
Attention computation and context vectors 这一部分把 End-to-End ASR: Attention 中相邻的教学单元串成了一条连续叙事：既交代了这一阶段要解决的问题，也交代了方法、推理或工程权衡为什么要这样组织。

\section{Training pathologies and alignment}
根据课程页主题、公开视频字幕以及本讲的 coverage ledger，Training pathologies and alignment 这一部分主要围绕 训练阶段的对齐失败、错误模式与 attention 何时会变得不稳定。 展开。
写作上保持 coverage-first 原则：先把老师在这一阶段真正讲过的问题设定、模型部件、训练或推理逻辑讲清，再压缩成适合复习的结构。

\subsection{训练阶段的对齐失败、错误模式与 attention 何时会变得不稳定}
从证据位置看，这一部分对应的课程证据区间是 00:14:51,560--00:22:33,880。课程在这里强调的核心内容可以概括为：训练阶段的对齐失败、错误模式与 attention 何时会变得不稳定。。
与单调模型不同，attention encoder-decoder 更强调全局条件化能力，但也因此引入了对齐自由度和部署层面的权衡。

在本章结构中，这一单元被并入“Training pathologies and alignment”，目的是把同一阶段的定义、机制和权衡集中讲清。

\subsection{本章小结}
Training pathologies and alignment 这一部分把 End-to-End ASR: Attention 中相邻的教学单元串成了一条连续叙事：既交代了这一阶段要解决的问题，也交代了方法、推理或工程权衡为什么要这样组织。

\section{Decoding and beam search}
根据课程页主题、公开视频字幕以及本讲的 coverage ledger，Decoding and beam search 这一部分主要围绕 beam search、解码路径与 token / frame 对齐在生成式 ASR 中的作用。 展开。
写作上保持 coverage-first 原则：先把老师在这一阶段真正讲过的问题设定、模型部件、训练或推理逻辑讲清，再压缩成适合复习的结构。

\subsection{beam search、解码路径与 token / frame 对齐在生成式 ASR 中的作用}
从证据位置看，这一部分对应的课程证据区间是 00:22:30,200--00:30:15,880。课程在这里强调的核心内容可以概括为：beam search、解码路径与 token / frame 对齐在生成式 ASR 中的作用。。
Search 这一讲把前面所有模型真正接到了 decoding pipeline 上，也就是把训练得到的分数转成最终转写结果。

在本章结构中，这一单元被并入“Decoding and beam search”，目的是把同一阶段的定义、机制和权衡集中讲清。

\subsection{本章小结}
Decoding and beam search 这一部分把 End-to-End ASR: Attention 中相邻的教学单元串成了一条连续叙事：既交代了这一阶段要解决的问题，也交代了方法、推理或工程权衡为什么要这样组织。

\section{Translation and LM connections}
根据课程页主题、公开视频字幕以及本讲的 coverage ledger，Translation and LM connections 这一部分主要围绕 translation / language-model 连接，以及 attention-based 方法与更大生成模型的关系。 展开。
写作上保持 coverage-first 原则：先把老师在这一阶段真正讲过的问题设定、模型部件、训练或推理逻辑讲清，再压缩成适合复习的结构。

\subsection{translation / language-model 连接，以及 attention-based 方法与更大生成模型的关系}
从证据位置看，这一部分对应的课程证据区间是 00:30:11,120--00:37:46,040。课程在这里强调的核心内容可以概括为：translation / language-model 连接，以及 attention-based 方法与更大生成模型的关系。。
与单调模型不同，attention encoder-decoder 更强调全局条件化能力，但也因此引入了对齐自由度和部署层面的权衡。

在本章结构中，这一单元被并入“Translation and LM connections”，目的是把同一阶段的定义、机制和权衡集中讲清。

\subsection{本章小结}
Translation and LM connections 这一部分把 End-to-End ASR: Attention 中相邻的教学单元串成了一条连续叙事：既交代了这一阶段要解决的问题，也交代了方法、推理或工程权衡为什么要这样组织。

\section{Transition to Transformer attention}
根据课程页主题、公开视频字幕以及本讲的 coverage ledger，Transition to Transformer attention 这一部分主要围绕 从 RNN attention 过渡到 Transformer attention，说明注意力如何成为主导范式。 展开。
写作上保持 coverage-first 原则：先把老师在这一阶段真正讲过的问题设定、模型部件、训练或推理逻辑讲清，再压缩成适合复习的结构。

\subsection{从 RNN attention 过渡到 Transformer attention，说明注意力如何成为主导范式}
从证据位置看，这一部分对应的课程证据区间是 00:37:43,920--00:44:45,720。课程在这里强调的核心内容可以概括为：从 RNN attention 过渡到 Transformer attention，说明注意力如何成为主导范式。。
与单调模型不同，attention encoder-decoder 更强调全局条件化能力，但也因此引入了对齐自由度和部署层面的权衡。

在本章结构中，这一单元被并入“Transition to Transformer attention”，目的是把同一阶段的定义、机制和权衡集中讲清。

\subsection{本章小结}
Transition to Transformer attention 这一部分把 End-to-End ASR: Attention 中相邻的教学单元串成了一条连续叙事：既交代了这一阶段要解决的问题，也交代了方法、推理或工程权衡为什么要这样组织。

\section{Stability and encoder design bridge}
根据课程页主题、公开视频字幕以及本讲的 coverage ledger，Stability and encoder design bridge 这一部分主要围绕 attention 计算、稳定训练与后续现代 encoder 设计之间的连接。 展开。
写作上保持 coverage-first 原则：先把老师在这一阶段真正讲过的问题设定、模型部件、训练或推理逻辑讲清，再压缩成适合复习的结构。

\subsection{attention 计算、稳定训练与后续现代 encoder 设计之间的连接}
从证据位置看，这一部分对应的课程证据区间是 00:44:43,720--00:51:51,880。课程在这里强调的核心内容可以概括为：attention 计算、稳定训练与后续现代 encoder 设计之间的连接。。
与单调模型不同，attention encoder-decoder 更强调全局条件化能力，但也因此引入了对齐自由度和部署层面的权衡。

在本章结构中，这一单元被并入“Stability and encoder design bridge”，目的是把同一阶段的定义、机制和权衡集中讲清。

\subsection{本章小结}
Stability and encoder design bridge 这一部分把 End-to-End ASR: Attention 中相邻的教学单元串成了一条连续叙事：既交代了这一阶段要解决的问题，也交代了方法、推理或工程权衡为什么要这样组织。

\section{Recap and transition}
根据课程页主题、公开视频字幕以及本讲的 coverage ledger，Recap and transition 这一部分主要围绕 整讲收束：attention-based ASR 的优势、局限，以及下一讲过渡。 展开。
写作上保持 coverage-first 原则：先把老师在这一阶段真正讲过的问题设定、模型部件、训练或推理逻辑讲清，再压缩成适合复习的结构。

\subsection{整讲收束}
从证据位置看，这一部分对应的课程证据区间是 00:51:49,240--00:58:27,839。课程在这里强调的核心内容可以概括为：整讲收束：attention-based ASR 的优势、局限，以及下一讲过渡。。
与单调模型不同，attention encoder-decoder 更强调全局条件化能力，但也因此引入了对齐自由度和部署层面的权衡。

在本章结构中，这一单元被并入“Recap and transition”，目的是把同一阶段的定义、机制和权衡集中讲清。

\subsection{本章小结}
Recap and transition 这一部分把 End-to-End ASR: Attention 中相邻的教学单元串成了一条连续叙事：既交代了这一阶段要解决的问题，也交代了方法、推理或工程权衡为什么要这样组织。

\section{总结与延伸}
End-to-End ASR: Attention 这一讲在整门《Speech Recognition and Understanding》课程中的位置，是把 E2E ASR: Attention 讲成一个可复用的建模模块。
从教材化角度看，本讲最重要的不是记住单一术语，而是理解它如何嵌入完整的 automatic speech recognition pipeline：输入语音如何被表示、模型如何被训练、解码如何得到最终转写，以及这些设计分别带来什么假设和权衡。
由于公开源以课程页和公开视频字幕为主、缺少官方 slide PDF，本章在正文里尽量把术语、方法和工程语境解释完整，并把缺失项留在 omission 和 manifest 里，而不是凭空补写。
\end{document}
